\section*{Chapitre \ref{ch:iv2:your-strategy-and-your-project} --- Votre stratégie et votre projet}
\markboth{\protect\omTocMark{}{Réponses~: Votre stratégie et votre projet}}{\protect\omTocMark{}{Réponses~: Votre stratégie et votre projet}}

\begin{solution}{iq:iv2:your-strategy-and-your-project:1}
Dire quel était le problème, ce que l'on a fait, ce qui en est résulté, en chiffres, et ce que l'on remettrait en question. Puis s'arrêter, en moins de deux minutes, et laisser le recruteur choisir la suite.

Pour construire la sienne, écrire une phrase sur chacun des points suivants~: le problème, l'approche, la part qu'on a prise soi-même, le résultat et la principale faiblesse. La chronométrer à voix haute. Vérifier chaque chiffre~: il faut pouvoir en refaire le calcul au tableau et dire comment on l'a testé.

\textbf{Illustration, avec ses propres éléments entre crochets.} «~Je voulais savoir [question], parce que [pourquoi cela comptait]. J'ai utilisé [données] et [méthode]. Ma part a été [ce que l'on a fait]. Le résultat était [nombre], face à une référence de [nombre]. Je m'y fie moins à cause de [faiblesse], et je testerais cela avec [étape suivante].~»

\iqwrong{Raconter les choses dans l'ordre où on les a faites. Cela prend dix minutes, et le résultat n'arrive qu'en dernier, si toutefois il arrive.}
\iqfollow{Pourquoi avoir choisi cette méthode plutôt qu'une méthode plus simple~?}{Dire ce que donnait la méthode plus simple. Si on ne l'a jamais lancée, le dire, et expliquer comment on les comparerait.}
\end{solution}

\begin{solution}{iq:iv2:your-strategy-and-your-project:2}
Environ 2 points d'alpha, au plus. La part due au marché était de $1.4 \times 20~\% = 28~\%$, ce qui laisse $30 - 28 = 2~\%$.

Et les 2~\% d'une seule année restent largement dans le bruit. La bonne réponse dit donc que cette année ne montre pas de compétence, et parle plutôt du processus.

\iqwrong{Comparer 30 à 20 et revendiquer 10 points. Si l'on détient 1.4 fois le marché, on a gagné le rendement du marché sur 1.4 fois le capital engagé.}
\end{solution}

\begin{solution}{iq:iv2:your-strategy-and-your-project:3}
L'erreur type vaut environ 1.6 points, donc 60~\% se trouvent à environ 6.3 erreurs types au-dessus du pile ou face. Oui, sur ce test, on fait mieux que le pile ou face. La question suivante est de savoir si l'ensemble de test était propre.

Avec $p = 0.5$ sous l'hypothèse nulle, $\sqrt{0.25/1000} = 1.58~\%$, et $10/1.58 = 6.3$.

Les questions suivantes sont donc~: combien de modèles a-t-on essayés sur cet ensemble de test, comment les cas ont-ils été découpés, et les 60~\% valent-ils encore quelque chose après coûts~? Un résultat significatif peut rester inutile.

\iqwrong{S'arrêter au score $z$. Si le même ensemble de test a servi à choisir le modèle, ou si des lignes voisines se sont glissées dans l'entraînement, le 6.3 ne veut rien dire.}
\end{solution}

\begin{solution}{iq:iv2:your-strategy-and-your-project:4}
Nommer une décision précise, dire ce que l'on sait maintenant et que l'on ignorait alors, et dire ce qu'elle aurait changé. Choisir une décision réelle qui a coûté quelque chose~: un recruteur ne croit ni à «~rien~» ni à «~plus de temps~».

Pour construire la sienne, lister les décisions qui ont pris le plus de temps ou sur lesquelles on est revenu. En choisir une qu'on sait expliquer avec son effet, idéalement avec un chiffre.

\textbf{Illustration, avec ses propres éléments entre crochets.} «~J'aurais d'abord [construit la référence ou les contrôles de données]. J'ai commencé par [méthode complexe], et je n'ai découvert [problème] qu'à [étape tardive]. Avec [contrôle simple] dès le départ, j'aurais gagné [temps] et trouvé [résultat] plus tôt.~»

\iqwrong{Une vantardise déguisée, du type «~J'aurais visé encore plus haut~». Cela signale au recruteur qu'on ne cherche pas ses propres erreurs.}
\end{solution}

\begin{solution}{iq:iv2:your-strategy-and-your-project:5}
Le nommer, et dire pourquoi~: l'échantillon était court, un choix a été fait après avoir vu les données, ou il repose sur une seule hypothèse. Puis dire quel test trancherait. Douter d'un de ses propres résultats est la meilleure preuve de jugement, et cela rend les autres plus faciles à croire.

Pour préparer la sienne, passer en revue chaque chiffre que l'on prévoit de citer, et se demander lequel on attaquerait si un concurrent l'avait rédigé.

\textbf{Illustration, avec ses propres éléments entre crochets.} «~Le [résultat] est le plus fragile. Il repose sur [échantillon court ou une seule hypothèse], et j'ai choisi [paramètre] après avoir vu [quoi]. Je le relancerais avec [nouvelles données ou un choix fixé] avant de m'y fier.~»

\iqwrong{Dire qu'on croit toutes ces réponses également. Personne d'expérimenté ne croit cela de son propre travail, et la question suivante débusquera la plus faible sans aide.}
\end{solution}

\begin{solution}{iq:iv2:your-strategy-and-your-project:6}
Oui, cela pourrait très bien être le cas. Un ratio de Sharpe véritable de zéro produit un historique aussi bon environ 3.3~\% du temps, donc environ 7 des 200 afficheraient un tel résultat sans aucune compétence.

La statistique $t$ vaut $1.5\sqrt{1.5} = 1.84$, et $\Phi(-1.84) = 3.3~\%$. Donc $200 \times 3.3~\% = 6.6$.

L'historique ouvre donc la discussion, pas le verdict. Ce qui aide, c'est de montrer le processus~: pourquoi l'on trade ce que l'on trade, et comment les transactions effectives se comparent au plan. Dire aussi combien d'autres pistes on a essayées puis abandonnées.

\iqwrong{Prendre 18 bons mois pour une preuve. Les recruteurs voient les candidats qui ont survécu, et l'on ne voit soi-même que les stratégies qui ont survécu.}
\end{solution}

\begin{solution}{iq:iv2:your-strategy-and-your-project:7}
Donner le résultat en une phrase qu'un non-spécialiste intelligent peut suivre, puis la méthode en trois. Ensuite, être honnête sur l'utilité~: nommer la compétence qui se transpose, même si le sujet, lui, ne se transpose pas.

Pour construire la sienne, rédiger la thèse en trois longueurs~: une phrase, une minute et dix minutes. Se préparer à la question d'un mathématicien sur l'argument central, et à la question «~et alors~?~».

\textbf{Illustration, avec ses propres éléments entre crochets.} «~J'ai montré que [résultat], ce qui compte parce que [raison simple]. Le sujet lui-même ne fera pas gagner d'argent ici. Ce qui se transpose, c'est [estimer à partir de données bruitées, ou démontrer qu'un résultat tient sous des hypothèses limitées], que j'utiliserais sur [problème précis].~»

\iqwrong{Forcer un lien avec le trading qui n'existe pas. Le recruteur le sait, et cela coûte plus cher que de dire que le sujet n'a pas d'utilité ici.}
\iqfollow{Quelle est l'étape la plus difficile de la démonstration, ou l'étape principale de l'estimation~?}{Être prêt à la refaire au tableau, sans notes. C'est la vérification que la démonstration est bien la sienne.}
\end{solution}

\begin{solution}{iq:iv2:your-strategy-and-your-project:8}
Nommer un type de contrepartie, et dire pourquoi elle négocierait à ce prix pour une raison autre que la valeur~: un vendeur contraint, un indice qui doit se rééquilibrer, un hedger qui paie pour une assurance, ou un acteur plus lent à réagir aux nouvelles. Si l'on ne sait pas qui c'est, le perdant le plus probable, c'est soi-même.

Pour préparer la sienne, compléter la phrase «~[ce groupe] négocie ceci parce que [règle ou contrainte], si bien que le prix s'écarte d'environ [ampleur], jusqu'à [événement].~» Puis se demander pourquoi un fonds plus gros n'a pas déjà pris la position.

\textbf{Illustration, avec ses propres éléments entre crochets.} «~De l'autre côté se trouve [groupe], qui doit [contrainte] avant [date], quel que soit le prix. Le prix reste donc à environ [ampleur] de la fair value. C'est trop petit pour [grands acteurs] parce que [raison]. Si la contrainte de [groupe] disparaît, ma raison d'entrer dans cette transaction disparaît aussi.~»

\iqwrong{«~Le marché a tort.~» Cela ne dit rien ni sur qui a tort, ni sur les raisons pour lesquelles il le resterait. Le recruteur entend une opinion, mais aucun avantage.}
\end{solution}

\begin{solution}{iq:iv2:your-strategy-and-your-project:9}
S'arrêter et le vérifier. Dire «~Laissez-moi regarder~», relire l'étape, et si le recruteur a raison, le reconnaître tout de suite et sans détour. Puis reconstruire à partir de ce point~: ce qui tombe, ce qui reste, et ce que l'on conclurait alors.

\begin{enumerate}
\item Ne pas défendre l'étape pendant qu'on la vérifie.
\item Dire en une phrase claire où l'on s'est trompé.
\item Reprendre à voix haute la partie touchée, et donner la nouvelle conclusion.
\item Si l'on reste en désaccord après vérification, dire pourquoi, avec un petit exemple numérique vérifiable.
\end{enumerate}

Pour construire cette habitude, s'entraîner avec quelqu'un qui interrompt à propos de vraies erreurs. Se faire corriger est souvent un piège tendu exprès~: on veut observer cette minute-là plus que la réponse.

\iqwrong{Discuter pendant plusieurs minutes, ou s'effondrer et abandonner la question. L'un comme l'autre montre comment on réagirait face à une position perdante, ce que le recruteur cherche justement à évaluer.}
\iqfollow{Et si vous pensez que c'est le recruteur qui a tort~?}{Le dire poliment, et le tester sur un petit cas numérique. Si l'on teste le candidat, tenir bon sur une étape vérifiée est justement ce qu'on attend.}
\end{solution}

\begin{solution}{iq:iv2:your-strategy-and-your-project:10}
Refuser poliment, et décrire ce qu'on peut dire en termes généraux. On ne partage pas le travail confidentiel d'un ancien employeur, et le dire montre qu'on ferait de même avec le travail de l'employeur qui recrute.

Pour préparer la sienne, décider avant l'entretien de ce qui est public ou générique~: le type de problème, le type de données, et la manière dont on l'a testé et fait tourner. Vérifier ce que dit son contrat, et ne pas deviner.

\textbf{Illustration, avec ses propres éléments entre crochets.} «~Je ne peux pas entrer dans les détails, car ils appartiennent à [employeur]. En termes généraux, j'ai travaillé sur [catégorie publique de signal], sur [horizon large], et j'avais la responsabilité de [ma part]. Je peux expliquer comment je l'aborderais à partir de zéro, avec des méthodes publiques.~»

\iqwrong{Faire allusion, avec un clin d'œil, aux signaux. Une société qui entend cela suppose qu'on ferait de même avec elle, et elle prend aussi un risque juridique en écoutant.}
\iqfollow{On veut juste un aperçu. Vous pouvez tout de même nous dire l'horizon et l'univers~?}{Seul ce qui a été vérifié est public ou générique. Réaffirmer son offre de traiter le problème à partir de zéro.}
\end{solution}

\begin{solution}{iq:iv2:your-strategy-and-your-project:11}
Dire exactement ce qu'on a fait, et nommer ce qu'ont fait les autres. On dit «~je~» pour sa part, et «~nous~» ou un rôle pour la leur. Si l'on ne peut pas reproduire le chiffre seul, dire sur quelle partie on aurait besoin d'aide.

Pour préparer la sienne, noter sur son CV sa part de chaque résultat. Si la part honnête est petite, commencer par ce qu'on a appris, et par la manière dont on le referait.

\textbf{Illustration, avec ses propres éléments entre crochets.} «~L'équipe a amélioré [métrique] de [chiffre]. J'ai construit [composant], ce qui représentait [ce qu'on peut honnêtement revendiquer, ou dire qu'on ne peut pas isoler]. [Rôle] a fait [autre partie]. Je pourrais reconstruire [composant] aujourd'hui, mais il me faudrait réapprendre [autre partie].~»

\iqwrong{«~Nous~» partout, ou «~je~» partout. Le premier laisse deviner ce qu'on a fait, et le second s'effondre au premier détail. Les recruteurs vérifient en descendant de deux niveaux.}
\end{solution}

\begin{solution}{iq:iv2:your-strategy-and-your-project:12}
L'espérance de rendement est de 2.8~\%. Le trade reste positif tant que le taux de réussite de 60~\% reste au-dessus d'environ 44.4~\%. Kelly donne un effet de levier de 3.5. Je ne l'utiliserais pas~: j'en prendrais nettement moins de la moitié, compte tenu de l'incertitude sur les 60~\%.

$0.6 \times 10 - 0.4 \times 8 = 2.8$. Le seuil d'équilibre est $8/18 = 44.4~\%$.

Pour un trade à deux issues, de gain $w$ et de perte $l$, $f = p/l - (1 - p)/w = 7.5 - 4 = 3.5$.

À 3.5 fois, la perte de 8~\% fait perdre 28~\% du capital. Et cela suppose que les issues se limitent à ces deux-là, et que la probabilité est juste. Avec $p = 44.4~\%$, Kelly donne zéro, donc 16 points d'erreur suffisent à faire passer la taille de 3.5 fois à zéro.

\iqwrong{Donner le chiffre de Kelly comme taille de la position. Il faut une probabilité juste, et celle que l'on estime est précisément ce que le recruteur cherche à tester.}
\end{solution}

\begin{solution}{iq:iv2:your-strategy-and-your-project:13}
Environ 14~\% des cas signalés sont des fraudes. Sur 1\,000 cas, on attend 10 fraudes, dont 8 sont détectées, et environ 50 bons cas signalés~: 58 signalements, dont 8 sont justes.

La précision vaut $0.008/(0.008 + 0.0495) = 13.9~\%$.

Donc une AUC de 0.95 dit peu de chose sur l'usage. Montrer une courbe précision-rappel au taux de base, ainsi que le coût de chaque type d'erreur. Le choix du seuil devient alors une décision métier, que l'on peut défendre.

\iqwrong{Commencer par l'AUC ou l'exactitude sur des événements rares. Répondre toujours «~aucune fraude~» donne 99~\% d'exactitude.}
\end{solution}

\begin{solution}{iq:iv2:your-strategy-and-your-project:14}
Environ 2 points sur les 80~\% relèvent de l'effet de sélection si les soumissions sont indépendantes, et environ 0.64 à une corrélation de 0.9. Pour distinguer son score de celui d'un rival situé 0.3 point derrière, il faudrait environ 137\,000 lignes avec des erreurs indépendantes, et environ 27\,000 avec des erreurs corrélées à 0.8.

\textbf{Sélection.} Un score public a une erreur type de $\sqrt{0.8 \times 0.2/2000} = 0.894$ point. Le meilleur de 50 tirages indépendants a une espérance de 2.25 erreurs types. Donc $2.25 \times 0.894 = 2.0$ points.

\textbf{Corrélation.} Avec une corrélation $\rho$ égale partout, seule la part indépendante se disperse, donc le facteur vaut $\sqrt{1 - \rho} = 0.316$, ce qui donne 0.64 point.

\textbf{Séparation.} Avec des erreurs indépendantes, l'écart a une erreur type de $\sqrt{2 \times 0.16/N}$. Pour que $0.003$ vaille 1.96 de ces erreurs types, $N = 2 \times 0.16 \times (1.96/0.003)^2 = 136\,600$. Sur les mêmes lignes, avec des erreurs corrélées entre modèles, la variance de l'écart diminue avec la corrélation~: à 0.8, elle vaut $2 \times 0.16 \times 0.2/N$, ce qui demande environ 27\,300 lignes.

Donc un classement à la 10\up{e} place sur 2\,000 tient surtout au processus de travail, et la façon honnête de le présenter est de donner le score final sur des données qu'on n'a pas vues.

\iqwrong{Présenter son score public, ou une place au classement, comme sa compétence. Demander ce que disait l'ensemble privé, et combien d'autres étaient à moins d'une erreur type.}
\end{solution}

\begin{solution}{iq:iv2:your-strategy-and-your-project:15}
On attend environ 19~\%, et seules environ 7~\% des trajectoires font aussi bien que 12~\%. Ce n'est donc pas une bonne nouvelle~: soit l'échantillon a été clément, soit l'historique présenté a été choisi, et il ne faut pas dimensionner sur 12~\%. Un jour ordinaire, on s'attend à reculer d'environ 5 à 6~\% par rapport au plus haut.

\textbf{Simulation.} Avec une marche aléatoire gaussienne quotidienne à ces paramètres, le maximum sur 5 ans a une moyenne d'environ 19~\%, une médiane d'environ 18~\%, et un 95\up{e}~centile d'environ 31~\%. Un drawdown supérieur à 30~\% a environ 6~\% de chances de se produire, même si la stratégie est exactement aussi bonne qu'on le dit.

\textbf{Formule stationnaire.} Pour un mouvement brownien avec dérive $\mu$ et volatilité $\sigma$, la distance au plus haut courant suit à long terme une loi exponentielle, de moyenne $\sigma^2/(2\mu) = 0.0225/0.4 = 5.6~\%$. Cela montre aussi qu'une dérive plus forte réduit vite les drawdowns, et qu'ils croissent avec le carré de la volatilité.

Donc, pour fixer un seuil de sortie, il faut utiliser un chiffre attendu, et non le chiffre favorable que l'on a obtenu. Avec cette simulation, une limite à 12~\% arrête une stratégie saine dans environ 93 cas sur 100 en cinq ans.

\iqwrong{Présenter le drawdown de son échantillon comme son risque. C'est un tirage dans une distribution large, et ici du côté bas.}
\iqfollow{Et si le rendement réel n'était que la moitié~?}{Le drawdown espéré monte à environ 26~\%, car la dérive le retient moins. Refaire la simulation avec la nouvelle dérive avant de fixer une limite.}
\end{solution}

\begin{solution}{iq:iv2:your-strategy-and-your-project:16}
Le test donne une moyenne d'environ 1 microseconde et un 99\up{e}~centile de 1 microseconde. Un flux à son propre rythme donne une moyenne d'environ 56 microsecondes, un 99\up{e}~centile d'environ 1 milliseconde, et un 99.9\up{e}~centile d'environ 9 millisecondes. Il aurait fallu mesurer à partir de l'instant d'envoi prévu de chaque message.

\textbf{Ce qui a été mesuré.} En une seconde, on envoie environ 990\,000 messages, à 1 microseconde chacun, et l'un prend 10 millisecondes. Un seul échantillon lent est donc noyé parmi 990\,000 bons, et la moyenne vaut 1.01 microseconde.

\textbf{Le flux.} Un message qui arrive $t$ millisecondes après le début du blocage attend la fin du blocage et les $0.1t$ de travail qui le précèdent~: $10 - 0.9t$ millisecondes. La file continue de se vider à ce rythme après la fin du blocage, donc l'attente décroît en ligne droite de 10 millisecondes à zéro sur $10/0.9 = 11.1$ millisecondes. La moyenne vaut $\tfrac12 \times 10 \times 11.1/1000 = 56$ microsecondes. Pour la queue de la distribution, la fraction de la seconde avec une attente supérieure à $w$ vaut $(10 - w)/0.9$ millièmes. En fixant cette fraction à 1~\%, on obtient $w = 1$ milliseconde, et à 0.1~\%, on obtient $w = 9.1$ millisecondes.

\textbf{Que faire.} Le problème s'appelle omission coordonnée (coordinated omission)~: le test cesse d'envoyer justement quand le système va mal. Il faut donc utiliser un générateur en boucle ouverte, à calendrier fixe, mesurer à partir de l'instant où chaque message aurait dû partir, et présenter la queue de distribution avec le débit. Puis trouver la cause du blocage.

\iqwrong{Présenter une moyenne, ou dire qu'un message lent sur une seconde n'a pas d'importance. Dans un système de trading, le blocage survient justement quand les cours bougent, et tous les messages arrivent à ce moment-là.}
\end{solution}

\begin{solution}{iq:iv2:your-strategy-and-your-project:17}
Le ratio de Sharpe vaut environ 0.95 avec ces jours, et environ 0.39 sans eux. Une distribution normale donnerait environ 14 points sur les 5 meilleurs jours, et non 40. Le résultat dépend donc de quelques événements, et la question est de savoir si l'on peut les expliquer.

\textbf{Ratios de Sharpe.} Avec eux, la moyenne vaut 6 points de base par jour pour une volatilité de 100, et $0.06 \times \sqrt{252} = 0.95$. Sans eux, la moyenne vaut $20/995 = 2.0$ points de base, et la volatilité baisse aussi, car chacun de ces jours vaut 8 points~: la variance perd $5 \times 64$ sur environ 1\,004 points carrés, si bien que la volatilité vaut environ 83 points de base. On obtient ainsi $(2.01/82.9) \times \sqrt{252} = 0.385$.

\textbf{Loi normale.} Retirer les 5 meilleurs jours de n'importe quel échantillon abaisse sa moyenne, il faut donc comparer avec le cas normal. Les 0.5~\% de valeurs les plus élevées d'une loi normale ont une moyenne de $\varphi(2.576)/0.005 = 2.89$ écarts types. Donc $5 \times 2.89 = 14.5$ points. Le nombre exact pour un échantillon fini est environ 14.3. Une stratégie normale symétrique perdrait tout de même 14 points sans eux, et celle-ci en perd 40.

\textbf{Ce que cela dit.} Chacun des meilleurs jours vaut environ 8 écarts types. Il s'agit soit d'une vraie source d'avantage (edge) que l'on peut nommer (par exemple une crise dont la stratégie est conçue pour tirer profit), soit d'une erreur de données, soit de la chance. Vérifier d'abord les cinq dates. Si elles sont réelles, montrer alors que ce type d'événement doit revenir, et à quelle fréquence. Si l'on retire les meilleurs jours, retirer aussi les pires, sinon la comparaison est biaisée.

\iqwrong{Dire que le résultat est robuste parce qu'il reste positif sans les meilleurs jours. C'est l'ampleur de la baisse, comparée à celle d'une loi normale, qui indique à quel point il en dépend.}
\end{solution}

\begin{solution}{iq:iv2:your-strategy-and-your-project:18}
La probabilité qu'il soit vrai est d'environ 72~\%. Au seuil de 5~\%, elle est d'environ 34~\%, et au seuil de 0.1~\%, d'environ 96~\%.

Sur 1\,000 idées, 50 sont vraies, et la moitié d'entre elles sont identifiées~: 25. Sur les 950 fausses, 1~\% passent~: 9.5. Donc $25/34.5 = 72~\%$.

Au seuil de 5~\%~: $25/(25 + 47.5) = 34~\%$. Au seuil de 0.1~\%~: $25/(25 + 0.95) = 96~\%$, en supposant la même puissance, ce qui n'est pas réaliste, car un seuil plus strict fait perdre de la puissance.

\textbf{Ce que je dis.} «~Le résultat est significatif à 1~\%, ce qui, selon mon estimation du taux de base dans le domaine, donne un résultat réel environ 7 fois sur 10. Pas 99~\%. C'est aussi un seul test, et je voudrais le valider hors échantillon (out-of-sample).~»

\iqwrong{Dire qu'une $p$-valeur de 1~\% signifie une chance de 99~\% que le résultat soit réel. Cela confond la probabilité des données sachant qu'il n'y a pas d'effet avec la probabilité qu'il n'y ait pas d'effet sachant les données.}
\end{solution}

\begin{solution}{iq:iv2:your-strategy-and-your-project:19}
La statistique de test vaut environ 1.86 si l'on suppose l'indépendance, et environ 1.37 en tenant compte de l'autocorrélation. En une phrase~: la perte moyenne de mon modèle était plus faible, mais avec 500 jours je ne peux pas exclure que ce soit du bruit.

$t = 0.05/0.6 \times \sqrt{500} = 1.86$. Pour un AR(1) de coefficient $\rho$, la variance de long terme est multipliée par $(1 + \rho)/(1 - \rho) = 1.857$, donc le $t$ est divisé par 1.363.

Il s'agit du test de Diebold--Mariano~: un test $t$ sur les différences de pertes, avec une variance de long terme. Et les deux prévisions doivent être comparées équitablement, aux mêmes dates, sans que l'une ni l'autre n'ait été ajustée sur la période de test.

\iqwrong{Dire «~il a battu l'indice de référence~» sur la seule moyenne, ou utiliser l'erreur type usuelle sur des données dont les erreurs sont persistantes.}
\end{solution}

\begin{solution}{iq:iv2:your-strategy-and-your-project:20}
(a) 3.0. (b) Un rendement d'environ 14.2~\% et un ratio de Sharpe d'environ 1.19. (c) Alpha de 14.4~\%, et ratio d'information d'environ 1.31. (d) Pas à lui seul~: 1.4 ne représente qu'environ 1.4 erreur type. (e) Environ 37. (f) Je retiens un Sharpe d'environ 1, et non 1.5, et je le lancerais à petite taille jusqu'à ce que les exécutions et la décroissance sur la seconde moitié soient expliquées.

\textbf{(a)} Le ratio de Sharpe vaut $18/12 = 1.5$, donc $t = 1.5\sqrt{4} = 3.0$.

\textbf{(b)} Le coût vaut $0.5 \times 3 \times 252 = 378$ points de base par an. Le doubler ramène le rendement à $18 - 3.78 = 14.22~\%$, et le ratio de Sharpe vaut 1.19.

\textbf{(c)} $18 - 0.3 \times 12 = 14.4$. La variance résiduelle vaut $0.0144 - 0.09 \times 0.0256 = 0.0121$, donc sa racine vaut 11~\%, et le ratio vaut $14.4/11 = 1.31$.

\textbf{(d)} Chaque moitié dure 2 ans, donc chacune a une erreur type de $1/\sqrt{2}$. Leur différence a une erreur type de 1. Donc $2.2 - 0.8 = 1.4$ vaut 1.4 erreur type, avec une $p$-valeur bilatérale de 16~\%. Ce n'est ni une preuve de décroissance, ni une preuve du contraire. Il faut regarder le signal et les coûts, qui sont plus parlants que le P\&L.

\textbf{(e)} Unilatéral, $\Phi(-3) = 0.135~\%$. Pour que la probabilité qu'au moins une variante parmi $N$ passe à tort atteigne 5~\%~: $N = \ln 0.95/\ln(1 - 0.00135) = 38$, ou $0.05/0.00135 = 37$ par la règle simple. Je dois donc pouvoir dire que j'en ai essayé bien moins que cela, et compter aussi les variantes corrélées.

\textbf{(f)} Dire quelque chose comme ceci. «~Trois éléments font baisser le chiffre~: la seconde moitié, un coût réaliste et ma recherche. Je n'ai pas de test qui les sépare en quatre ans. Je retiens donc un Sharpe d'environ 1, et le résultat net des coûts est ce que je vérifierais en premier en réel, car le coût par dollar se mesure en quelques semaines. Je le dimensionnerais en conséquence, et non sur 1.5. J'arrêterais si le coefficient d'information (IC) mesuré du signal baissait, ou si les coûts réels dépassaient le double de ceux du modèle. Ce sont des éléments vérifiables. Un mauvais trimestre, non.~»

\iqwrong{Répondre correctement à chaque question sans jamais arriver à une conclusion. La dernière question est celle qui intéresse vraiment le recruteur, et il faut donner un chiffre et un plan.}
\iqfollow{Vous me dites que le Sharpe de 1 relève d'un jugement. Comment l'avez-vous obtenu~?}{Reconnaître franchement qu'il s'agit d'un jugement~: le 1.5 réduit pour tenir compte du réalisme des coûts, et une décote de sélection qui dépend de la recherche. Montrer la fourchette, et dire ce qui ferait changer d'avis.}
\end{solution}
